% Generated draft: requires mathematical and source review.
\documentclass[11pt]{article}
\usepackage[margin=25mm]{geometry}
\usepackage{fontspec}
\setmainfont{Latin Modern Roman}
\usepackage{amsmath,amssymb,mathtools,unicode-math}
\setmathfont{Latin Modern Math}
\usepackage{enumitem,needspace,xurl,hyperref}
\setlength{\parindent}{0pt}
\setlength{\parskip}{5pt}
\setlength{\emergencystretch}{3em}
\newcommand{\N}{\mathbb{N}}
\newcommand{\Z}{\mathbb{Z}}
\newcommand{\Q}{\mathbb{Q}}
\newcommand{\R}{\mathbb{R}}
\newcommand{\C}{\mathbb{C}}
\newcommand{\F}{\mathbb{F}}
\newcommand{\A}{\mathbb{A}}
\newcommand{\PP}{\mathbb{P}}
\newcommand{\E}{\mathbb{E}}
\newcommand{\eps}{\varepsilon}
\newcommand{\dd}{\,\mathrm{d}}
\begin{document}
\section*{UnsolvedMath ID: 11100001}
\textbf{Source problem code:} catalogue entry
\subsection{Definitions and mathematical statement}
Mark Hovey's first major item is a request to formulate a specific, coherent and somewhat accessible programme for homotopy theory, analogous in organizing role to the Weil conjectures in algebraic geometry and the Ravenel conjectures in the late 1970s. It is not a stated mathematical proposition with objects, hypotheses and a conclusion to prove.

A response in the source's intended spirit would articulate concrete conjectures and explain how their relations supply a shared direction for the field. The source gives no required list, formal adequacy criterion or unique target theorem, so no quantified existence assertion or new conjectural equivalence is substituted. ``Analogous'' means a comparable research-organizing role, not logical equivalence to the Weil conjectures.

Hovey says this was his motivation for making the list: a forum in which conjectures could collectively form such a programme. He distinguishes the talents of formulating conjectures from proving them, illustrates the point with Deligne/Weil and Hopkins/Ravenel, and uses Fields-Medal language rhetorically. Those personal illustrations are motivation, not mathematical conditions or a prize criterion. He also explicitly cautions that the proposed conjectures should be somewhat accessible and that the subsequent listed problems might not be.

The research summary names later chromatic-homotopy directions, Lurie's higher/derived-algebraic-geometric programme with derived moduli of formal groups or elliptic curves, and Stolz--Teichner links between elliptic cohomology and two-dimensional quantum field theory as related examples. They are reported illustrations, not a verified uniquely accepted answer or a supplied set of conjectures resolving this subjective request.
\subsection{Short English statement}
Hovey asks the community to formulate an accessible, unifying programme of concrete homotopy-theory conjectures. This is a forum-style research aspiration, with no objective theorem or unique solution specified.
\subsection{Sources}
\begin{itemize}
\item Local UnsolvedMath registry, permanent ID \(11100001\), exact source code \texttt{catalogue entry}.
\item Mark Hovey, \emph{Hovey's Algebraic Topology Problem List}, Major problems, item 1, surviving mirror: \url{https://www-users.cse.umn.edu/~tlawson/hovey/}. Registry access 2026-07-29; html-ordered-list extraction.
\item Weil, Ravenel, Deligne and Mike Hopkins appear as the source's motivational examples. The summary's chromatic, Lurie and Stolz--Teichner directions have no supplied precise bibliographic citations.
\end{itemize}
\end{document}
